\documentclass[11pt,a4paper]{article}
\usepackage[T1]{fontenc}
\usepackage{lmodern,amsmath,amssymb,amsthm,microtype}
\usepackage[margin=26mm]{geometry}
\usepackage[colorlinks=true,linkcolor=blue,urlcolor=blue]{hyperref}
\newtheorem{theorem}{Theorem}[section]
\newtheorem{lemma}[theorem]{Lemma}
\newtheorem{corollary}[theorem]{Corollary}
\newcommand{\Q}{\mathbb Q}\newcommand{\Z}{\mathbb Z}
\title{Input-only weights do not complete an elliptic index comparison}
\author{}\date{17 November 2026}
\begin{document}\maketitle

\begin{abstract}
An additional balanced pole with a polynomial weight in the input
can force a genuine upper bound for the input height. It does not
identify that input with the integer quotient of two elliptic
indices. For a specific degree-eight weight, every input $2+1/q$
outside a fixed finite set of denominator primes has a complete
rational witness. The construction uses a local inverse for the
truncated-logarithm ratio and the Chinese remainder theorem.
All fixed-place, real, denominator and finite-character guards are
retained. The obstruction concerns the stated source, not arbitrary
sources with two elliptic points.
\end{abstract}

\section{Local interpolation of an index ratio}

Fix $d\in\Z_{>0}$ and $E_d:y^2=x^3-d^2x$, with formal
parameter $z=-x/y$. Let $\Lambda(Z)\in\Q[Z]$ satisfy
$\Lambda(0)=0$, $\Lambda'(0)=1$. At every prime $p$, choose a
depth $e_p$ such that the formal logarithm identifies
$E_{e_p}(\Q_p)$ with $p^{e_p}\Z_p$ isometrically and
\[
 p^{-e}\Lambda(p^e Z)=Z+pG_e(Z),\qquad G_e\in\Z_p[Z],
                         \qquad e\ge e_p.
\]
Such a depth exists by the formal-group logarithm
\cite[Chapters IV and VII]{AEC}. If $p$ is odd and the coefficients of
$\Lambda$ are $p$-integral, one may take $e_p=1$.

\begin{lemma}\label{lem:interpolation}
Let $H\in E_{e_p}(\Q_p)$ be nonzero. The continuous map
\[
 \Phi_H:\Z_p\longrightarrow\Z_p,\qquad
 a\longmapsto\frac{\Lambda(z([a]H))}{\Lambda(z(H))},
\]
where $[a]H$ is defined through the formal logarithm, is an
isometric bijection. In particular, for every $X\in\Z_p$ and
every $N\ge1$, there is a residue class of integers $a$ modulo
$p^N$ for which $v_p(\Phi_H(a)-X)\ge N$.
\end{lemma}
\begin{proof}
Write $e=v_p(z(H))$. The normalized polynomial above is an
isometric bijection of $\Z_p$, by Hensel's lemma. Thus
$z\mapsto\Lambda(z)$ is an isometric bijection of
$p^e\Z_p$. The formal logarithm and its inverse are also
isometries there. Consequently the four maps
\[
 a\mapsto a\log_E(H),\quad
 a\log_E(H)\mapsto z([a]H),\quad
 z\mapsto\Lambda(z),\quad
 t\mapsto t/\Lambda(z(H))
\]
compose to an isometric bijection of $\Z_p$. The two scaling
maps contribute opposite valuation shifts $e$ and $-e$.
The last assertion follows by reducing this bijection modulo
$p^N$.
\end{proof}

\section{The two pole blocks}

Fix a rank-one congruent-number curve and a nontorsion point
$R\in2E_d(\Q)$, and put $P=2R$. Choose $R$ sufficiently deep in the formal
groups at two and seven. Orient $R$ so that
\[
 q_0=\frac{x(R)^2+d^2}{2y(R)}>0,\qquad q_0^2=x(P).
\]
Choose a positive rational $\lambda_0$ with
$v_2(\lambda_0)=2$, $v_7(\lambda_0)=-1$ and
$\lambda_0^2>100$, and put $c=q_0\lambda_0^2$.

Fix $T\ge1$, let $\Lambda=\Lambda_T$ be the degree-at-most-$T$
truncation of the formal logarithm, and put $k=8T-7$. Retain
points $H=[h]P$, $U=[j]P$, rational square abscissas
$q_H^2=x(H)$ and $q_U^2=x(U)$, and the comparison test
\[
             C=\frac{\Lambda(z(U))}{\Lambda(z(H))}-X.
\]
The first block has weight and tests
\[
        u_1=(c/q_H)^k,\qquad C_{1,1}=x(U)^{-T},
                          \qquad C_{1,2}=C.
\]
The second block has
\begin{equation}\label{eq:input-weight}
 \begin{split}
 F(X)&=\prod_{a\in\{1,3,5,9\}}((56X)^2-a^2),\\
 u_2&=16F(X),\qquad C_{2,1}=x(H)^{-T},
                         \qquad C_{2,2}=C.
 \end{split}
\end{equation}
For $i=1,2$ define
\[
 D_{i,r}=\frac{C_{i,r}^2}{1+u_iC_{i,r}^2},\qquad
 W_i=\prod_{r=1}^2(1+D_{i,r}^2/u_i),\qquad
 b_i=\frac{4u_iW_i^2}{(1+u_i)^2},\qquad \eta_i=7b_i.
\]
Impose the two intrinsic equations
\begin{equation}\label{eq:intrinsic}
 Y_i^2=\eta_i(\eta_i+v_i^2)^2s_i^4
       -\eta_i(\eta_i-v_i^2)^2(1-\eta_i)^2.
\end{equation}
Keep the point, quotient, weight and intrinsic norm guards
of the index-pole source \cite{Index}. Exclude any fixed finite
set of indices $h$. The all-index primitive-denominator theorem
under weak $abc$ with exponent $4/3$ may be used; no coprime-index
progression is imposed here.

\begin{theorem}\label{thm:false-inputs}
There is a fixed finite set $S_0$ of primes such that every prime
$q\notin S_0$ gives a noninteger input
\[
                              X=2+1/q
\]
having complete rational witnesses for the two-block source.
The witnesses may be taken with $U=[m]H$ for an integer
$m\ne0$, and with arbitrarily large $|h|$.
Every finite-character guard is nonzero at the prepared rational
parameters. In particular, the second pole does not make this
source a definition of $\Z$.
\end{theorem}

\section{Preparation of the reference point}

For a point $T$ near the identity write
\[
                q(T)=\frac{x(T)^2+d^2}{2y(T)},
                       \qquad q(T)^2=x([2]T).
\]
The product $q(T)\log_E(T)$ is analytic at the identity with
nonzero limiting value $-1/2$. By choosing $R$ sufficiently
deep at two and seven, uniformly for every nonzero integer $h$,
\begin{equation}\label{eq:square-ratio}
 \frac{q_0}{q([h]R)}=h\epsilon_{p,h},\qquad
 \epsilon_{2,h}\in1+8\Z_2,\qquad
 \epsilon_{7,h}\in1+7\Z_7.
\end{equation}
Indeed $\log_E([h]R)=h\log_E(R)$ and every $[h]R$ stays in
the same fixed formal neighborhood.
Both displayed unit groups consist of local squares.

Fix $q$ outside the bad reduction, coefficient, logarithm and
constant primes, also excluding two and seven. Put $X=2+1/q$.
The weight $u_2$ is positive and larger than $50$,
\[
 u_2\in\Q_2^{\times2},\quad v_2(u_2)=4,\qquad
 u_2\in\Q_7^{\times2},\quad v_7(u_2)=0.
\]
These facts follow from $56X\in8\Z_2\cap7\Z_7$, the four
odd integers $1,3,5,9$, and $X>2$. At $q$ one has
\[
                            v_q(u_2)=-8.
\]
At every other prime $X$ is integral and $u_2$ is integral.

Choose $H=[h]P$ formal, to any prescribed sufficiently large
depth, at all fixed exceptional primes, at $q$, and at every
positive odd prime of $u_2$. At $q$ require
$v_q(z(H))$ to be even. These conditions can be imposed on
an arithmetic progression of integers $h$ that are local squares
at two and seven. To see this, choose an integer $L$ such that
$[L]P$ is sufficiently deep at all these finitely many places.
Write $h=L^2r$. Impose $r\equiv1\pmod8$ and
$r\equiv1\pmod7$. At $q$, choose $v_q(r)$ to be zero or one,
as necessary to make the denominator depth of $H$ even, and
fix the corresponding nonzero residue modulo the next power
of $q$. These congruences are compatible. Increasing $L$
beforehand supplies every other required depth.

For completeness of this assertion, on a sufficiently deep
formal subgroup,
$v_p(z([a]Q))=v_p(z(Q))+v_p(a)$ for integral nonzero $a$.
Thus the selected parity at $q$ is exact. The required finite
depths can be chosen so that, whenever $v_p(u_i)>0$ is odd,
the first detector factor is a unit:
\[
 4v_p(C_{1,1})-v_p(u_1)>0,\qquad
 4v_p(C_{2,1})-v_p(u_2)>0.
\]
For the first inequality, outside the constant primes it is
at least $(8T-k)v_p(z(H))=7v_p(z(H))>0$.
At a constant prime it follows after increasing the chosen
depth. The second inequality has a fixed right-hand threshold
once $X$ is fixed.

Take $[h]R$ in a sufficiently small real neighborhood of $R$,
so that the coherent root $q_H=q([h]R)$ is positive and
$c/q_H>50$. This is possible for infinitely many arbitrarily
large $h$ in the chosen arithmetic progression: multiples of
a nontorsion point are dense in the appropriate real component.
Here $R\in2E_d(\Q)$ lies in the identity real component.
The fixed affine and logarithm guards remove finitely many
indices. Formula~\eqref{eq:square-ratio} gives
\begin{equation}\label{eq:weight-fixed}
 \begin{split}
 &u_1\in\Q_2^{\times2},\qquad
       v_2(u_1)=k(4+v_2(h))\ge4k,\\
 &u_1\in\Q_7^{\times2},\qquad
       v_7(u_1)=k(-2+v_7(h)).
 \end{split}
\end{equation}
At $q$, the valuation of $u_1$ is even by the prescribed
denominator parity. This is the only prime where $X$ is not
integral.

\section{Choosing the second index and the intrinsic witnesses}

After fixing $H$, let $\mathcal T$ contain two, seven, all fixed
exceptional primes, and every positive odd prime of $u_1$ or
$u_2$, omitting $q$. At each prime of $\mathcal T$ where a
comparison is required, $H$ is formal and $X\in\Z_p$.
At primes arising only from the denominator of $H$, this follows
automatically; at the other primes it was imposed above.
The interpolation lemma applies, with the enlarged fixed depths
at coefficient primes.

Choose an integer $m\ne0$ by the Chinese remainder theorem so
that, with $U=[m]H$,
\[
 v_p(C)>N_p\qquad(p\in\mathcal T),
\]
where $N_p$ exceeds one quarter of every positive $v_p(u_i)$
present at that prime. All such choices form an arithmetic
progression, so finitely many additional index exclusions may
be avoided.

At every positive odd weight valuation, both detector factors
are now units, and $v_p(b_i)=v_p(u_i)>0$. At a negative odd
valuation $\nu=v_p(u_i)$, the balanced identity
\[
             v_p(D_{i,r})=\max(2v_p(C_{i,r}),-\nu)
\]
makes every factor of $W_i$ a unit and gives
$v_p(b_i)=-\nu>0$. At an even weight valuation the valuation
of $b_i$ is even, regardless of cancellation. In particular,
the denominator prime $q$ is harmless for both blocks.

At two put $e_P=v_2(z(P))\ge4$. Then
$e_H=e_P+v_2(h)$ and
\[
 4v_2(C_{1,1})-v_2(u_1)
   \ge8T(e_P+v_2(h))-k(4+v_2(h))
   \ge28+7v_2(h)>0.
\]
The first factor of the second block is also a unit, since
$8Te_H>v_2(u_2)=4$. Both comparison factors are units by
interpolation. Therefore
\[
 b_i\in\Q_2^{\times2},\qquad
                     v_2(b_i)=2+v_2(u_i)\ge6.
\]
At seven the square class of $u_i$ prevents cancellation
against $-1$. If $v_7(u_i)\le0$, both detector factors are
units automatically. If $v_7(u_i)>0$, the depth and interpolation
conditions make them units. Since $1+u_i$ is a unit when
$v_7(u_i)=0$, one obtains $v_7(b_i)\ge0$ in every case.
Thus both $\eta_i=7b_i$ have square reduced denominator and
satisfy the required two-adic class and depth.

At every real point character, $u_i>50$ and all tests are real.
Consequently
\[
 0\le D_{i,r}<1/u_i,\qquad
 1\le W_i\le(1+u_i^{-3})^2,\qquad
                  0<b_i<1/7.
\]
The balanced denominators, products and inputs are therefore
nonzero at every character, independently of the size of $m$.
Finite isogeny covers do not change these image values. The
intrinsic realization theorem \cite{Intrinsic} supplies rational
pairs $(v_i,s_i)$ for \eqref{eq:intrinsic}. Its Zariski density
allows simultaneous avoidance of the radicand norms and every
retained guard whose restriction to the corresponding intrinsic
solution variety is nonzero. This proves
Theorem~\ref{thm:false-inputs}.

\section{The input-height estimate still holds}

Under ordinary rational $abc$, the polynomial-form radical
estimate of \cite{Derivative} gives, outside a fixed set $S$,
\[
 \log r^+_{\mathrm{odd},S}(u_2)\ge3h(X)-O(1).
\]
Indeed the homogeneous form
\[
           \prod_{a\in\{1,3,5,9\}}(56^2A^2-a^2B^2)
\]
is squarefree of degree eight. Its radical has logarithm at
least $(6-\epsilon)h(A:B)-O(1)$, while its logarithmic size
is at most $8h(A:B)+O(1)$. The simple-prime support therefore
has logarithm at least $(4-2\epsilon)h(A:B)-O(1)$.
Use $\epsilon=1/2$. Rational cancellation occurs only at the
fixed coefficient primes.

The second pole forces $H$ formal at every consumed odd prime,
since its first test is $x(H)^{-T}$. Hence
\[
                          3h(X)\le\widehat h(H)+O(1).
\]
For $T=4$ and the weak exponent $4/3$, the amplified first
block also implies, for every false input $X\ne m$,
\[
              \log(1+|m|)\ge\tfrac53\widehat h(H)-O(1).
\]
The witnesses above have room for this growth because the
second index is chosen after all finite local constraints.
An input-height bound alone therefore leaves the index-size
term uncontrolled.

\begin{thebibliography}{9}
\bibitem{Index}\emph{Elliptic index divisibility and amplified congruences},
companion manuscript,
\nolinkurl{elliptic_index_poles_20261117_v60.tex}.
\bibitem{Intrinsic}\emph{A two-coordinate pole test of coefficient
height four}, companion manuscript,
\nolinkurl{intrinsic_pole_20261117_v56.tex}.
\bibitem{Derivative}\emph{An eleven-variable construction under
rational $abc$}, companion manuscript,
\nolinkurl{derivative_detector11_20261117_v59.tex}, Section 1.
\bibitem{AEC}J.~H.~Silverman, \emph{The Arithmetic of Elliptic Curves},
second edition, Graduate Texts in Mathematics 106, Springer, 2009,
Chapters IV and VII.
\end{thebibliography}
\end{document}
